% from ega1-6-fr.tex
\oldpage[I]{153}
restriction de $f$ à $U$ soit un morphisme de type fini de $U$ dans $V$. On
dit alors aussi que $X$ est un préschéma localement de type fini sur $Y$, ou
un $Y$-préschéma localement de type fini.
\end{definition}

Il résulte aussitôt de (\hyperref[I.6.3.2-fr]{6.3.2}) que si $f$ est
localement de type fini, alors, pour tout ouvert $W$ de $Y$, la restriction
de $f$ à $f^{-1}(W)$ est un morphisme $f^{-1}(W)\to W$ qui est localement de
type fini.

Si $Y$ est localement noethérien et si $X$ est localement de type fini sur
$Y$, $X$ est localement noethérien en vertu de
(\hyperref[I.6.3.7-fr]{6.3.7}).

\begin{proposition}[6.6.3]
\label{I.6.6.3-fr}
Pour qu'un morphisme $f:X\to Y$ soit de type fini, il faut et il suffit qu'il
soit quasi-compact et localement de type fini.
\end{proposition}

La nécessité des conditions est immédiate, vu
(\hyperref[I.6.3.1-fr]{6.3.1}) et la remarque suivant
(\hyperref[I.6.6.1-fr]{6.6.1}). Inversement, supposons ces conditions
satisfaites et soit $U$ un ouvert affine de $Y$, d'anneau $A$~; pour tout
$x\in f^{-1}(U)$, il y a par hypothèse un voisinage
$V(x)\subset f^{-1}(U)$ de $x$ et un voisinage $W(x)\subset U$ de
$y=f(x)$, contenant $f(V(x))$ et tel que la restriction de $f$ à $V(x)$ soit
un morphisme $V(x)\to W(x)$ qui est de type fini. En remplaçant $W(x)$ par
un voisinage $W_1(x)\subset W(x)$ de $x$ de la forme $D(g)$ (avec $g\in A$),
et $V(x)$ par $V(x)\cap f^{-1}(W_1(x))$, on peut supposer que $W(x)$ est de
la forme $D(g)$, donc de type fini sur $U$ (puisque son anneau s'écrit
$A[1/g]$)~; par suite $V(x)$ est de type fini sur $U$. En outre
$f^{-1}(U)$ est quasi-compact par hypothèse, donc réunion d'un nombre fini
d'ouverts $V(x_i)$, ce qui achève la démonstration.

\begin{proposition}[6.6.4]
\label{I.6.6.4-fr}
\begin{enumerate}
  \item[(i)] Une immersion $X\to Y$ est quasi-compacte si elle est fermée,
    ou si l'espace sous-jacent à $Y$ est localement noethérien ou si l'espace
    sous-jacent à $X$ est noethérien.
  \item[(ii)] Le composé de deux morphismes quasi-compacts est quasi-compact.
  \item[(iii)] Si $f:X\to Y$ est un $S$-morphisme quasi-compact, il en est de
    même de $f_{(S')}:X_{(S')}\to Y_{(S')}$ pour toute extension $g:S'\to S$
    du préschéma de base.
  \item[(iv)] Si $f:X\to X'$ et $g:Y\to Y'$ sont deux $S$-morphismes
    quasi-compacts, $f\times_S g$ est quasi-compact.
  \item[(v)] Si le composé de deux morphismes $f:X\to Y$, $g:Y\to Z$ est
    quasi-compact et si $g$ est séparé, ou l'espace sous-jacent à $X$
    localement noethérien, $f$ est quasi-compact.
  \item[(vi)] Pour qu'un morphisme $f$ soit quasi-compact, il faut et il
    suffit que $f_{\mathrm{red}}$ le soit.
\end{enumerate}
\end{proposition}

On notera que (vi) est évident puisque la propriété d'être quasi-compact pour
un morphisme ne dépend que de l'application continue correspondante des
espaces sous-jacents. Démontrons de même la partie de (v) correspondant au cas
où l'espace sous-jacent $X$ est supposé localement noethérien. Posons
$h=g\circ f$, et soit $U$ un ouvert quasi-compact dans $Y$~; $g(U)$ est
quasi-compact (non nécessairement ouvert) dans $Z$, donc contenu dans une
réunion finie d'ouverts quasi-compacts $V_j$
(\hyperref[I.2.1.3-fr]{2.1.3}), et $f^{-1}(U)$ est par suite contenu dans la
réunion des $h^{-1}(V_j)$, qui sont des sous-espaces quasi-compacts de $X$,
donc des sous-espaces noethériens. On en conclut
(\hyperref[0.2.2.3-fr]{0, 2.2.3}) que $f^{-1}(U)$ est un espace noethérien,
et \emph{a fortiori} quasi-compact.

Pour prouver les autres assertions, il suffit de démontrer (i), (ii) et (iii)
(\hyperref[I.5.5.12-fr]{5.5.12}). Or, (ii) est évidente, et (i) résulte de
(\hyperref[I.6.3.5-fr]{6.3.5}) lorsque l'espace $Y$ est localement noethérien
ou l'espace $X$ noethérien et est évidente pour une immersion fermée. Pour
établir (iii),
